\section{Connected pieces of the midpoint fan}
\label{sec:area_law_fan_pieces}

A vector $(u_1,u_2)\in\R^2$ has an allowed direction if
$u_1=0$, $u_2=0$, $u_1=u_2$, or $u_1=-u_2$.
We apply this condition to the radial segments of the elementary fans in
\cite[Section~11]{OpenAI2026AreaLaw}.

\begin{theorem}[Allowed directions in an affine span]
    \label{thm:area_law_allowed_direction_affine_span}
    \lean{TNLean.PEPS.AreaLaw.Geometry.isAllowedSlope_sub_of_mem_affineSpan_pair}
    \leanok
    Let $p,q\in\R^2$ and suppose that $q-p$ has an allowed direction.
    If $v,x\in\operatorname{aff}\{p,q\}$, then $x-v$ has an allowed
    direction.
\end{theorem}
\begin{proof}\leanok
    Write $v=(1-r)p+rq$ and $x=(1-s)p+sq$ with $r,s\in\R$.
    Then $x-v=(s-r)(q-p)$, and each of the four linear equations defining
    an allowed direction is preserved under scalar multiplication.
\end{proof}

\begin{theorem}[Allowed base and radial directions]
    \label{thm:area_law_fan_base_radial_directions}
    \lean{TNLean.PEPS.AreaLaw.Geometry.cellFanPolygon_base_and_radial_isAllowedSlope}
    \leanok
    \uses{def:area_law_cell_fan}
    In an elementary fan, with each side of the cell either undivided or
    subdivided at its midpoint, the vectors $b_i-a_i$ and $c-b_i$ have
    allowed directions for every fan triangle $P_i=\operatorname{conv}\{c,a_i,b_i\}$.
\end{theorem}
\begin{proof}\leanok
    \uses{def:area_law_template,def:area_law_cell_fan}
    Each fan triangle with vertices $c,a_i,b_i$ is an allowed polygonal piece,
    so its edge vectors $b_i-a_i$ and $c-b_i$ have slopes in
    $\{0,\infty,1,-1\}$.
\end{proof}

Now subdivide every side of the dyadic square at its midpoint, giving eight
triangles $P_i$ and eight radial segments $R_i=[c,b_i]$.

\begin{theorem}[Nonempty connected pieces after radial deletion]
    \label{thm:area_law_fan_cut_piece_connected}
    \lean{TNLean.PEPS.AreaLaw.Geometry.cellFanPolygon_sdiff_radials_inter_ball_nonempty}
    \lean{TNLean.PEPS.AreaLaw.Geometry.cellFanPolygon_sdiff_radials_inter_ball_isConnected}
    \leanok
    \uses{def:area_law_cell_fan}
    Let $A$ be any subset of the eight radial indices and let $r>0$.
    For each triangle index $i$, the set
    \begin{align}
        \left(P_i\setminus\bigcup_{t\in A}R_t\right)\cap B_\infty(c,r)
    \end{align}
    is nonempty and connected.
\end{theorem}
\begin{proof}\leanok
    \uses{thm:area_law_cell_fan_regularity,
        thm:area_law_cell_fan_sector_directions,
        thm:area_law_allowed_direction_affine_span,
        thm:area_law_fan_base_radial_directions,
        thm:area_law_fan_radial_union_deletion_convex}
    Since $P_i=\overline{\operatorname{int}(P_i)}$ and $c\in P_i$, the ball
    $B_\infty(c,r)$ contains a point $x\in\operatorname{int}(P_i)$.
    The displacement $x-c$ has none of the four allowed directions, whereas
    $y-c$ has an allowed direction for every point $y$ on any radial $R_t$.
    Thus $x\notin R_t$ for all $t$, so the intersection is nonempty.

    The triangle with the selected radials deleted is convex, as is the ball
    $B_\infty(c,r)$. Their intersection is therefore a nonempty convex set,
    hence connected.
\end{proof}
